% titel: Blatt 2, A5: Summenformel aus Massenanteilen und Molmasse
% skript: Mortimer Kap. 4.2 (empirische Formel, Molmasse)

% Teil I: Werkzeuge, kein Lösungsweg.

Aus den Massenanteilen (Prozentangaben) einer Verbindung laesst sich das
Verhaeltnis der Atome bestimmen. Kennt man zusaetzlich die experimentelle
Molmasse, findet man die wirkliche (Summen-)Formel --- nicht nur das
kleinste ganzzahlige Verhaeltnis.

\begin{werkzeugkasten}[Massanteil --- Stoffmenge-Verhaeltnis --- Summenformel]
Massenanteile $w_i$ (z.~B. $47{,}3\% = 0{,}473$) geben an, welcher Anteil der
Gesamtmasse auf jedes Element entfaellt. Mit der bekannten Molmasse $M$
berechnet man die Masse jedes Elements in einem mol:
\begin{equation}
  m_i = w_i \cdot M\qquad\rightarrow\qquad
  n_i = \frac{m_i}{M(\mathrm{Element}_i)}.
\end{equation}
Die Stoffmengenverhaeltnisse $n_i$ stehen im Verhaeltnis der Atomszahlen. Man
dividiert alle $n_i$ durch den kleinsten Wert und erhaelt ganzzahlige (bis auf
Rundung) Faktoren --- das ist zunaechst die \textbf{empirische Formel}. Die
experimentelle Molmasse $M$ entspricht entweder genau der Masse der empirischen
Formel oder einem ganzzahligen Vielfachen davon; so ergibt sich die
Summenformel.
\textit{Skript:} Mortimer Kap. 4.2 (empirische und molekulare Formel),
Übung 2/5.
\end{werkzeugkasten}

\begin{achtung}[Darauf achten]
\begin{itemize}
  \item Prozente vor dem Rechnen als Dezimalbruch einsetzen
        ($50{,}0\% = 0{,}500$) --- eine der haeufigsten Fehlerquellen.
  \item Massenanteile auf 355 g Gesamtmasse (die experimentelle Molmasse)
        beziehen; so arbeitet man direkt in $\mathrm{g/mol}$.
  \item Zum Schluss auf ganze Zahlen runden; Abweichungen von ein bis zwei
        Prozent sind normale Rundungsfehler.
  \item Die Summe der Massanteile ist $100\%$ --- zur Kontrolle prüfen
        ($47{,}3+2{,}54+50{,}0 = 99{,}8$).
\end{itemize}
\end{achtung}
